\section{Implementation and Module-Wise Coding}

\subsection{Development Environment}
\begin{itemize}
\item \textbf{Languages:} Python 3.13 (back end, about 10{,}500 lines in \texttt{backend/app}), TypeScript
      in strict mode (app, about 5{,}500 lines in \texttt{mobile/src}), SQL (migrations, templates).
\item \textbf{Frameworks and libraries:} FastAPI, SQLAlchemy 2.1, Alembic, Pydantic v2, PyMuPDF,
      pytesseract, openpyxl, python-docx, RapidFuzz, argon2-cffi, PyJWT; Expo SDK 58, React Native 0.88,
      Expo Router, TanStack Query, React Hook Form, Zod.
\item \textbf{Tools:} VS Code, Git/GitHub, PostgreSQL 16, Docker, Expo Go, Chrome DevTools, pytest,
      Jest, Playwright, Render.
\item \textbf{Operating systems:} Windows and Linux laptops for development; Android and iPhone for device testing.
\end{itemize}

\subsection{Overall Implementation Strategy}
The design was translated into software \emph{risk first}, from the inside out:
\begin{enumerate}
\item \textbf{Foundation:} schema, migrations, authentication, RBAC and admin master data, so every
      later feature was secure from the start.
\item \textbf{Extraction engine:} the riskiest part, built and tuned on the real college PDF early.
\item \textbf{Query engine:} the intent registry and templates, tested intent by intent (Q01--Q30).
\item \textbf{Mobile app} on top of a stable, documented API (OpenAPI).
\item \textbf{Hardening:} concurrency tests, a security audit, the real-PDF rules from the college,
      laptop/web support, the SDK upgrade and free cloud deployment.
\end{enumerate}
Each step ended with all tests passing and a commit (Table~\ref{tab:commits}).

\subsection{Module-Wise Implementation}
\paragraph{M1. Authentication and RBAC (\texttt{app/security}, \texttt{services/auth.py}).}
\emph{Input/Output:} credentials $\rightarrow$ access token (15 min) and refresh token (30 days, stored only as a hash).
\emph{Logic:} Argon2id verification; on refresh, the old token is revoked and a new one is issued in
the same family, and reuse of a revoked token revokes the family. \texttt{load\_principal} reads roles,
staff profile and academic placement from the database on every request. \texttt{policy.py} maps
roles to data scopes (batch, department, college).
\emph{Tests:} 15 tests in \texttt{test\_auth.py} (rotation, reuse detection, suspension, role change, rate limit).

\paragraph{M2. Upload and archive (\texttt{services/uploads.py}, \texttt{storage.py}).}
\emph{Input:} multipart file plus metadata. \emph{Logic:} extension allow-list and magic-byte
check (\texttt{extraction/filetypes.py}), size and page limits, atomic private write (mode 0600),
SHA-256 for duplicate detection, archive row and job row in one transaction, idempotency via
\texttt{client\_request\_id}. \emph{Output:} timetable id in state \textsc{Queued}.
\emph{Tests:} \texttt{test\_uploads\_and\_domains.py} (13).

\paragraph{M3. Extraction (\texttt{app/extraction}, 11 files).}
\texttt{pdf.py} (PyMuPDF tables and words), \texttt{geometric.py} (gutter-based columns),
\texttt{ocr.py} (Tesseract), \texttt{office.py} (XLSX/CSV/DOCX, LibreOffice for legacy formats)
$\rightarrow$ \texttt{model.py} (Grid) $\rightarrow$ \texttt{sections.py}, \texttt{timeparse.py},
\texttt{gridparse.py}, \texttt{cellparse.py} $\rightarrow$ \texttt{validate.py} $\rightarrow$
\texttt{pipeline.py}. The parser version (\texttt{tt-grid-parser/1.2.0}) is stored with every entry.
\emph{Tests:} 24 extraction tests and 14 regression tests on the real PDF.

\paragraph{M4. Worker (\texttt{app/worker.py}).}
Claims jobs (Listing~\ref{lst:claim}), runs the pipeline, persists entries, computes the timetable
status, notifies the uploader, and retries unexpected failures. It can run as a separate process
(\texttt{python -m app.worker}) or as a thread inside the API on the free tier.

\paragraph{M5. Query engine (\texttt{app/query}).} Files \texttt{nlp.py}, \texttt{vocab.py}, \texttt{registry.py}, \texttt{engine.py} and \texttt{templates.py}.
Implements Algorithm~\ref{alg:query}. \texttt{registry.py} declares the 30 intents with their allowed
domains, roles, required parameters and templates. \texttt{templates.py} holds the 22 parameterised
SQLAlchemy queries. \emph{Tests:} 57 tests in \texttt{test\_query\_engine.py}, including SQL-injection text
and unknown parameters.

\paragraph{M6. Conflicts (\texttt{services/conflicts.py}).}
Implements Algorithm~\ref{alg:conflict} and returns \textsc{Confirmed}/\textsc{Possible} conflicts,
scoped to the HOD's department. \emph{Test:} the real-PDF test asserts that exactly the two genuine clashes
are reported.

\paragraph{M7. Academics and administration.} Files \texttt{services/academics.py}, \texttt{users.py} and \texttt{api/v1/admin.py}.
Fresher registration without duplicates, promotion with preview and confirm (all or nothing, history
preserved), and master data CRUD with validation. \emph{Tests:} 13 in \texttt{test\_academics\_admin.py}, including
concurrent promotions.

\paragraph{M8. Mobile and web client (\texttt{mobile/src}).}
Described in Section~7.4. \emph{Tests:} 18 Jest tests and Playwright end-to-end tests.

\subsection{Front-End Implementation}
The app has 26 screens organised with Expo Router (file-based routing):
\begin{itemize}
\item \textbf{Auth stack:} welcome, login, account activation, forgot password.
\item \textbf{Tabs:} Home (role dashboard), Ask (NL search), Timetable (week strip and day timeline),
      Archives, Admin (Admins only; hiding the tab is cosmetic because the server enforces access), Profile.
\item \textbf{Stack screens:} class details, timetable details (extraction summary, sections),
      review \& correct, upload, conflicts, monitoring, notifications, admin tools (users, new user,
      promotion, master data, institution settings, jobs, audit log).
\end{itemize}
\textbf{Forms} use React Hook Form with Zod schemas (e.g.\ e-mail format, password rules, room code).
\textbf{Navigation} is file-based. \textbf{State}: server data in TanStack Query with keys per domain
and timetable, invalidated when the official timetable changes. A
\textbf{College/Personal toggle} switches the domain, with a different accent colour so users always
know which data they see. \textbf{Accessibility}: light/dark themes, large tap targets, labels for
screen readers. On a wide laptop screen the app is shown as a centred column. The web session is kept in
\texttt{sessionStorage} (cleared when the tab closes), and on phones in the Keychain/Keystore.

\subsection{Back-End Implementation}
\begin{itemize}
\item \textbf{76 REST endpoints} in six routers: \texttt{auth}, \texttt{me}, \texttt{timetables},
      \texttt{search}, \texttt{notifications}, \texttt{admin}. The OpenAPI specification is generated
      automatically (\texttt{docs/openapi.json}).
\item \textbf{Request processing:} middleware adds a request id and JSON logging; dependencies
      authenticate and load the principal; Pydantic validates bodies; services hold business logic; the
      error handler converts exceptions to the envelope.
\item \textbf{Search context:} \texttt{services/context.py} decides the domain (request hint
      $\rightarrow$ remembered $\rightarrow$ institutional), the timetable (explicit archive
      $\rightarrow$ remembered archive $\rightarrow$ primary), verifies ownership and usability, and
      computes ``now'' in the user's time zone.
\item \textbf{Static web app:} in the deployed image the same FastAPI process serves the exported web
      app, so phone and laptop use one HTTPS origin.
\end{itemize}

\subsection{Database Implementation}
31 tables are created by Alembic migrations. Templates are SQLAlchemy Core/ORM expressions with
bound parameters. A representative template (QT08, next class) is:
\begin{lstlisting}[language=Python,caption={Template QT08 (simplified): next class within a scope},label=lst:qt08]
stmt = (select(E)
        .where(E.timetable_id == ctx.timetable_id,
               E.entry_kind == "CLASS",
               E.verification_status != "REJECTED",
               E.time_uncertain.is_(False),
               or_(E.class_date == day, and_(E.class_date.is_(None),
                                             E.day_of_week == day.isoweekday())),
               E.start_time > after,
               scope_filter(ctx))          # batch + parent batch, or staff id
        .order_by(E.start_time)
        .limit(1))
\end{lstlisting}
Connectivity: one SQLAlchemy engine per process with a connection pool (\texttt{psycopg} 3 driver);
each request gets a session that commits or rolls back as one transaction.

\subsection{Communication and Concurrency Implementation}
\begin{itemize}
\item \textbf{Client--server:} \texttt{mobile/src/api/client.ts} adds the bearer token, applies
      timeouts, retries idempotent GETs, and performs a \emph{single-flight} refresh: the first 401
      starts one refresh call and all other failing requests wait for the same promise, then retry.
\item \textbf{Concurrent clients:} several Uvicorn processes; each request is independent
      (stateless), with its own database transaction.
\item \textbf{Synchronisation:} activation and preference changes use \texttt{SELECT \ldots{} FOR
      UPDATE}; job claiming uses \texttt{SKIP LOCKED}; uploads are idempotent through a partial unique index;
      the rate limiter uses a lock-protected in-process sliding window.
\item \textbf{Recovery:} job heartbeat and re-claim; maximum attempts; the client shows the cached
      institutional timetable with ``last refreshed'' time when offline.
\end{itemize}

\subsection{Git Repository Setup}
\begin{itemize}
\item \textbf{Repository URL:} \url{\RepoURL}
\item \textbf{Structure:} \texttt{backend/} (app, alembic, tests), \texttt{mobile/} (src, \_\_tests\_\_, e2e),
      \texttt{docs/} (13 documents + OpenAPI), \texttt{deploy/} (Render image), \texttt{scripts/}
      (dev script, Windows one-click start, load smoke), \texttt{report/} (this report and poster),
      \texttt{render.yaml}, \texttt{docker-compose.yml}.
\item \textbf{Branching strategy:} \texttt{main} is the stable branch; development happens on a feature
      branch that is merged through a pull request after review. The Render demo deploys from the
      development branch, so every push updates the demo.
\item \textbf{Commit message convention:} an imperative summary of what changed and why, naming the
      scope (e.g.\ ``Query engine: intent registry Q01--Q30 \ldots''), with the details in the body.
\end{itemize}

\begingroup\small
\begin{longtable}{@{}p{1.5cm}p{1.9cm}p{9.0cm}p{2.3cm}@{}}
\caption{Development history (commits on the development branch)}\label{tab:commits}\\
\toprule
\textbf{Commit} & \textbf{Date} & \textbf{Message (summary)} & \textbf{Files / lines}\\
\midrule
\endfirsthead
\toprule
\textbf{Commit} & \textbf{Date} & \textbf{Message (summary)} & \textbf{Files / lines}\\
\midrule
\endhead
\texttt{0dfda92} & 2026-10-09 & Backend foundation: schema, migrations, auth, RBAC, admin master data & 46 / +5164\\
\texttt{bc917fb} & 2026-10-09 & Extraction engine (PDF/OCR/Excel/CSV/Word/images), worker, archives, review, conflicts & 30 / +3841\\
\texttt{47e3d89} & 2026-10-09 & Query engine: intent registry Q01--Q30, templates QT01--QT22, entity resolution, clarification & 17 / +3026\\
\texttt{1aa651d} & 2026-10-09 & Expo mobile app: role dashboards, NL search, weekly timetable, archives, uploads, review, admin; Jest + Playwright & 72 / +18618\\
\texttt{f028feb} & 2026-10-09 & Academic/admin/conflict/concurrency tests, batched vocabulary loading, load smoke script & 7 / +353\\
\texttt{4f3528b} & 2026-10-09 & Documentation, OpenAPI export, Docker/compose, EAS profiles, pinned requirements & 23 / +7345\\
\texttt{bb81091} & 2026-10-09 & Tune extraction on the college's real timetable PDF (parser 1.1.0) & 13 / +508\\
\texttt{dc12b25} & 2026-10-09 & Apply the college's timetable rules; real-PDF seed; phone guide & 31 / +1106\\
\texttt{710e679} & 2026-10-14 & College seed: one student login per lab group & 3 / +24\\
\texttt{fce3644} & 2026-10-21 & Run the app in a laptop browser; one-click Windows start & 7 / +129\\
\texttt{e8f74c3} & 2026-10-21 & Hide admin review notices from students; fix archive banner & 12 / +102\\
\texttt{326b8e6} & 2026-10-25 & Upgrade the app to Expo SDK 58 & 8 / +9018 (lockfile)\\
\texttt{8093f09} & 2026-10-25 & Free online demo (Render) + fix account takeover via password reset (P0-1) & 25 / +444\\
\bottomrule
\end{longtable}
\endgroup

\paragraph{Contribution of individual group members.}
\begin{table}[H]
\centering
\caption{Contribution of group members (to be confirmed by the group and guide)}
\label{tab:contrib}
\small
\begin{tabularx}{\textwidth}{@{}p{3.6cm}X@{}}
\toprule
\textbf{Member} & \textbf{Main contributions}\\
\midrule
\StudentA{} (\UIDA) & \textit{[e.g.\ back-end foundation, query engine, deployment]}\\
\StudentB{} (\UIDB) & \textit{[e.g.\ extraction pipeline and real-PDF tuning]}\\
\StudentC{} (\UIDC) & \textit{[e.g.\ mobile app screens and UI design]}\\
\StudentD{} (\UIDD) & \textit{[e.g.\ testing, documentation, report and poster]}\\
\bottomrule
\end{tabularx}
\end{table}

\subsection{Integration of Modules}
Modules communicate only through defined interfaces: the app uses the REST API; routers call
services; services call the query engine or put jobs in the queue; the worker calls the extraction
pipeline and writes through the same models. Integration was verified at three levels:
(1) back-end integration tests that upload a file, run the worker, activate the timetable and query it
through the API; (2) the college seed test, which runs the real PDF through the whole pipeline and then
asks questions as students, faculty, HOD and Admin; (3) Playwright end-to-end tests that drive the real
UI against the real API, including uploading a PDF through the file picker and waiting for the worker.

\subsection{Implementation Challenges and Resolutions}
\begin{longtable}{@{}p{4.2cm}p{4.6cm}p{6.2cm}@{}}
\caption{Implementation challenges and resolutions}\label{tab:challenges}\\
\toprule
\textbf{Challenge} & \textbf{Cause} & \textbf{Resolution}\\
\midrule
\endfirsthead
\toprule
\textbf{Challenge} & \textbf{Cause} & \textbf{Resolution}\\
\midrule
\endhead
Noon row printed as ``12.15 a.m.\ to 01.15 p.m.'' & Typing error in the source PDF & Accept midday only when proven by both neighbouring rows (Algorithm~\ref{alg:time})\\
Stacked lab cells (A1--A4 in one cell) & Four sub-groups share one grid cell & Line grouping by sub-group token; each becomes its own entry with its own room/teacher\\
Legends split across columns, codes missing from legend & Page layout & Page-local legend parser; unknown codes kept as written, with an INFO note (not an error)\\
Combined B.E.\ electives across divisions A--D & One page serves four divisions & Map to a common parent batch so all B.E.\ divisions see it\\
LLC slots with no teacher or room & Co-curricular course & Treat as a valid class without teacher/room; app hides empty fields\\
542 raw overlaps & Parallel labs and electives are intended & Intended-overlap rules $\rightarrow$ 2 genuine clashes\\
Slow searches (about 50 SQL statements) & N+1 lazy loading of related rows & Eager loading and batched vocabulary $\rightarrow$ about 25\\
Web app stuck on ``Loading\ldots'' in production & Connectivity probe used \texttt{HEAD /}, server returned 405 & Accept \texttt{GET} and \texttt{HEAD} on web routes; regression test\\
Account takeover risk (P0-1) & Password-reset code could be returned/logged in non-dev configuration & Never return or log codes in production; 9 verification checks\\
Expo Go on iPhone required SDK 58 & Store Expo Go supports only the newest SDK & Upgraded to SDK 58 / RN 0.88; fixed type changes; mocked icon fonts in Jest\\
Free hosting allows one service & Render free tier & In-process worker thread using the same \texttt{SKIP LOCKED} claim; static web app served by FastAPI\\
PowerShell blocked \texttt{npx} on Windows & Execution policy & Documented \texttt{npx.cmd} / policy change; one-click \texttt{start-windows.bat}\\
\bottomrule
\end{longtable}

\newpage
